\chapter{Diseño experimental}\label{cap:diseno}

\section{Dataset FEVER}

\subsection{Composición y particiones}

El dataset FEVER \citep{thorne2018fever} contiene 185\,445 \emph{claims}
sintéticos generados por anotadores humanos a partir de
aproximadamente 50\,000 artículos de Wikipedia. Cada \emph{claim} está
etiquetado como \textsc{supported} (\textsc{sup}), \textsc{refuted}
(\textsc{ref}) o \textsc{not-enough-info} (\textsc{nei}). En el caso
de las dos primeras etiquetas, los anotadores aportan además al menos
un conjunto de evidencias en forma de \emph{(página, número de frase)}.

La distribución oficial de FEVER incluye particiones \emph{train},
\emph{dev} y \emph{test}; esta última no incorpora etiquetas
públicamente disponibles. El presente trabajo utiliza la partición
\emph{train} para ajuste de parámetros, una partición de validación
construida estratificando un 10\,\% del \emph{train} (semilla fija,
\texttt{random\_state=42}) y la partición oficial \emph{dev} como
conjunto de prueba para evaluación final.

\subsection{Preprocesado}

Para los experimentos se construye una tabla
$(\text{claim}, \text{evidencia}, \text{etiqueta})$ donde cada fila
corresponde a un par claim--evidencia. En el caso \textsc{nei}, sin
evidencia anotada, se muestrea aleatoriamente una frase del mismo
artículo de Wikipedia (estrategia \emph{nearest neighbor sampling}). La
implementación del preprocesado se encuentra en
\path{research-stats/research_stats/data/preprocess.py}.

\subsection{Distribución de clases}

La distribución conjunta de etiquetas en \emph{train} (tras
preprocesado) es aproximadamente equilibrada:
% TODO(resultados): completar con números reales tras descarga de FEVER.
\begin{table}[h]
  \centering
  \begin{tabular}{lrr}
    \toprule
    Etiqueta & Frecuencia & Proporción \\
    \midrule
    \textsc{supported}        & --- & --- \\
    \textsc{refuted}          & --- & --- \\
    \textsc{not-enough-info}  & --- & --- \\
    \midrule
    Total & --- & 1.00 \\
    \bottomrule
  \end{tabular}
  \caption{Distribución empírica de etiquetas en \emph{train}.}
  \label{tab:fever-dist}
\end{table}

\section{Métricas de evaluación}

Sea $C \in \R^{K \times K}$ la matriz de confusión con
$C_{ij}$ el número de ejemplos de clase verdadera $i$ predichos como
$j$. Definimos para cada clase $k$:
\[
  \mathrm{prec}_k = \frac{C_{kk}}{\sum_j C_{jk}}, \quad
  \mathrm{rec}_k = \frac{C_{kk}}{\sum_j C_{kj}}, \quad
  F1_k = \frac{2\,\mathrm{prec}_k \mathrm{rec}_k}{\mathrm{prec}_k + \mathrm{rec}_k}.
\]
La \emph{exactitud} (\emph{accuracy}) es
$\mathrm{acc} = N^{-1} \sum_k C_{kk}$ y el F1 \emph{macro} es
$F1_{\text{macro}} = K^{-1} \sum_k F1_k$.

Adicionalmente reportamos:
\begin{itemize}[leftmargin=*,itemsep=0.2em]
  \item \emph{Log-loss} multiclase (NLL):
        $\mathrm{NLL} = -N^{-1} \sum_i \log \hat{p}_{y_i}(x_i)$.
  \item \emph{Brier score} multiclase:
        $\mathrm{BS} = N^{-1} \sum_i \sum_k
        \bigl(\hat{p}_k(x_i) - \mathbb{1}\{y_i = k\}\bigr)^2$.
  \item \emph{Expected Calibration Error} (ECE) con $M = 15$ bins
        equidistantes en $[0, 1]$.
\end{itemize}

\section{Protocolo de evaluación}

Cada experimento se ejecuta con tres semillas (\texttt{42, 1337, 2025})
y se reporta la media y la desviación típica de las métricas. Para la
comparación entre el baseline y el modelo neuronal se aplica el test
de McNemar pareado (capítulo~\ref{cap:tests}) sobre la partición de
prueba; para la diferencia de F1 macro se construye un intervalo de
confianza por \emph{bootstrap} no paramétrico con $B = 10\,000$
remuestreos.

\section{Recursos computacionales}

Los modelos lineales se entrenan en CPU. El \emph{fine-tuning} de
DeBERTa-v3-base se realiza, según disponibilidad, en una GPU NVIDIA
de memoria $\geq$ 16\,GB con precisión mixta \texttt{fp16}. Las
configuraciones detalladas se reportan en el anexo.
